\documentclass[10pt,a4paper]{amsart}
\usepackage[margin=19mm]{geometry}
\usepackage{amsmath,amssymb,mathtools}
\usepackage[scr=rsfs,cal=euler]{mathalfa}
\usepackage{xfrac}
\usepackage[unicode,hidelinks,bookmarks=false]{hyperref}
\newcommand{\ie}{i.e.\ }
\newcommand{\cf}{cf.\ }
\newcommand{\resp}{respectively\ }
\newcommand{\Subsection}{\subsection}
\providecommand{\nobreakdash}{\nobreak\hbox{-}}
\setlength{\emergencystretch}{2em}
\multlinegap0pt
\usepackage{bm}
\usepackage{bbm}
\usepackage{dsfont}
\usepackage{xspace}
\usepackage{cancel}
\usepackage{mathtools}
\usepackage{cleveref}
\usepackage[shortlabels]{enumitem}
\usepackage{amsthm}
\makeatletter
\@ifundefined{theorem}{\newtheorem{theorem}{Theorem}}{}
\@ifundefined{corollary}{\newtheorem{corollary}[theorem]{Corollary}}{}
\@ifundefined{proposition}{\newtheorem{proposition}[theorem]{Proposition}}{}
\makeatother
\newcommand{\IA}{\mathbb{A}}
\newcommand{\IN}{\mathbb{N}}
\newcommand{\into}{\hookrightarrow}
\newcommand{\supp}{\mathrm{supp \,}}
\renewcommand{\H}{\mathrm{H}}
\renewcommand{\deg}{\mathrm{deg}}
\renewcommand{\dim}{\mathrm{dim}}
\renewcommand{\div}{\mathrm{div}}
\title[The Tate conjecture for surfaces of geometric genus one -- e]{The Tate conjecture for surfaces of geometric genus one -- embracing singularities}
\author{Haoyang Guo, Ziquan Yang}
\date{Source: arXiv:2501.18541v2 (CC BY 4.0)}
\begin{document}
\maketitle

\noindent\emph{Source and licence.} The Tate conjecture for surfaces of geometric genus one -- embracing singularities. Version arXiv:2501.18541v2 (CC BY 4.0), distributed under the Creative Commons Attribution 4.0 International licence, as recorded in the arXiv licence metadata for this exact version and shown on the abstract page https://arxiv.org/abs/2501.18541v2. Full rights notice: licenses/arxiv-2501.18541-CC-BY-4.0.txt.\par\smallskip

\noindent\textit{Editorial scope.} Counted unit: the Corollary whose source label is \texttt{cor: bound inseparable} in the article (source0.tex lines 1570--1573, source label \texttt{cor: bound inseparable}), delivered together with the article's own complete original proof of it (source0.tex lines 1574--1590, one \texttt{proof} environment of 17 lines). No mathematics is written, repaired or re-derived: statement and proof are the article's own text line for line. Required context retained uncounted: eqn: sum to zero (lines 1521--1524); prop: inspearable j (lines 1527--1538). Every definition, display and label of those lines is retained unchanged. Omitted from this package and not redistributed: 1--1520, 1525--1526, 1539--1569, 1591--3246. Nothing inside a retained excerpt is omitted or altered: every retained body is delivered exactly as the article's lines are written, so no omission receipt of any kind accompanies this package. Citations: the retained excerpts carry no citation command at all, so no bibliography entry of the article is transcribed and no reference is redistributed. Reference adaptation: none was needed and none was made; every reference inside the delivered unit resolves inside it, so no unresolved reference or citation is produced. Printed slips kept literal: no printed word, symbol, hypothesis or number of the counted statement or of its proof was altered. Layout and class: the article's own class is \texttt{amsart} with packages including geometry, amsmath, mathrsfs, amsthm, amsfonts, amssymb, latexsym, tikz, esint, mathtools, stmaryrd, biblatex, fontenc, hyperref; the wrapper is the lane's pinned \texttt{amsart} editorial wrapper at 10pt on A4, which restates the theorem-like environments the retained text needs and adds the article's own macro definitions that the retained text needs (\texttt{\textbackslash H}, \texttt{\textbackslash IA}, \texttt{\textbackslash IN}, \texttt{\textbackslash deg}, \texttt{\textbackslash dim}, \texttt{\textbackslash div}, \texttt{\textbackslash into}, \texttt{\textbackslash supp}), with extra wrapper packages (bm, bbm, dsfont, xspace, cancel, mathtools, cleveref, enumitem). The delivered numbering is the unit's own. Assets: no retained span contains a figure environment, a graphics-inclusion command, a PDF-page inclusion, a table or a TikZ picture; no image file is copied into this package, the package declares no asset dependency and no image file is redistributed with it. Licence: version arXiv:2501.18541v2 of this article is distributed under the Creative Commons Attribution 4.0 International licence, as recorded in the arXiv licence metadata for this exact version and shown on the abstract page https://arxiv.org/abs/2501.18541v2. A full rights notice is delivered at licenses/arxiv-2501.18541-CC-BY-4.0.txt. Characters: every retained excerpt is pure ASCII, so the delivered engine, XeLaTeX with its default Latin Modern text font, reports no missing character.
\begin{equation}
    \label{eqn: sum to zero}
    \oplus_{x \in \supp(\mathrm{div}(s))} m_x(\div(s)) \cdot x = e.
\end{equation}

\begin{proposition}
\label{prop: inspearable j}
    If the $j$-invariant map is nonconstant and inseparable, then $p = 5$ or $7$. When $p = 5$, there are two possibilities: 
    \begin{enumerate}[label=\upshape{(\roman*)}]
        \item $\deg j = 10$, and there are points $y, z \neq x$ in $C$, such that $\mathrm{div}(a_4) = \{4x \}$, $\mathrm{val}_x(a_6) = 1$, and $\mathrm{div}(\Delta) = \{2x, 5y,  5z\}$. 
        \item $\deg j = 5$, and there are points $x \neq y$ in $C$, such that $\div(a_4) = \{ 3x, y \}$, and $\div(a_6) = \{ 2x, 4y \}$.  
    \end{enumerate}
    When $p = 7$, there is one possibility: 
    \begin{enumerate}
        \item[\upshape{(iii)}] $\deg j = 7$, and there are points $x \neq y$ in $C$, such that $\div(a_4) = \{ 3x, y \}$, and $\div(a_6) = \{ x, 5y \}$. 
    \end{enumerate}
\end{proposition}

\begin{corollary}
\label{cor: bound inseparable}
    View the vector space $\H^0(C, L^4) \times \H^0(C, L^6)$ as the affine space $\IA^4 \times \IA^6 \simeq \IA^{10}$ over $k$. The set $I$ of $k$-points $(a_4, a_6)$ such that the corresponding elliptic surface $X$ has finite inseparable $j$-invariant map is contained in a subscheme of dimension at most $2$. 
\end{corollary}

\begin{proof}
     
    Let $\{ \Delta = 0 \}$ be the subscheme of $\IA^{12}$ such that $\Delta(a_4, a_6)$ vanishes identically, and set $\IA' := \IA^{12} \smallsetminus \{ \Delta = 0 \}$. 
    For every $r \in \IN$, let $C^{[r]}$ denote the Hilbert scheme of $r$-points on $C$, which is well known to be the same as the $r$th symmetric power as $C$ is a smooth curve. 
    Then the association $(a_4, a_6) \mapsto \mathrm{div}(\Delta)$ defines a morphism $\IA' \to C^{[12]}$. 
    Note that a point on $C^{[r]}$ can also be viewed as a degree $r$ effective divisor. 

    Let us start with case (i) in \cref{prop: inspearable j}. As $a_4 \neq 0 \in L^4(-4 x)$, $x$ is a $4$-torsion point under the group law by (\ref{eqn: sum to zero}). Hence there are at most $16$ different choices for $x$. Let us fix one. Since $h^0(L(-4x)) = 1$, there is a line in $\IA^4$ parametrizing those $a_4$ such that $\mathrm{div}(a_4) = 4x$. 
    Let $\IA''$ be the fiber of $\IA'$ over this line.  
    Let $C' \in C^{[2]}$ be the subscheme parametrizing those pairs $y + z$ such that $2x \oplus 5y \oplus 5z = e$.
    Let $Z \subseteq C^{[12]}$ be the image of the morphism $C' \to C^{[12]}$ given by $y + z \mapsto 2 x + 5y + 5z$. 
    Since two sections $s, s' \in \H^0(C, L^{12})$ define the same divisor if and only if $s \in k^\times s'$, every fiber of the morphism $\IA' \to C^{[12]}$ (and hence the composition $\IA'' \into \IA' \to C^{[12]}$) has dimension at most $1$.
    Now, $I$ is contained the set of $k$-points of $\IA'' \times_{C^{[12]}} C'$. 
    As $\dim C' = 1$, $\dim \IA'' \times_{C^{[12]}} C' \le 1 + 1 = 2$. 

    Now consider case (ii). The equations $3 x \oplus y = e$ and $2x \oplus 4y = e$ imply that $x$ is a $10$-torsion point. Hence there are finitely many of them. Fixing each $x$, each $a_4$ and $a_6$ are well defined up to scalars, so we obtain the dimension bound. Case (iii) is entirely similar. 
\end{proof}

\end{document}
